% !TEX root = task-1-2-1.tex
% Editable task bank. Append new BankTask blocks at the end of the document.
\documentclass[12pt,a4paper]{article}
\usepackage{physics-textbook}
\usepackage{graphicx,needspace}
\usetikzlibrary{decorations.pathmorphing}
\newif\ifShowSolutions
\ifdefined\TutorEdition
  \ShowSolutionstrue
\else
  \ShowSolutionsfalse % Change to \ShowSolutionstrue to display solutions.
\fi
\newcounter{banktask}
\newsavebox{\BankTaskBox}
\newcommand{\BankTaskContent}[3]{%
  \noindent\textbf{\thebanktask.}\quad #2\par
  \ifShowSolutions
    \Needspace{50mm}% Keep the solution heading with its opening text or diagram.
    #3\par
    \ifdefined\TutorEdition
      \ifcsname TutorSource#1\endcsname
        \nopagebreak\smallskip{\small\csname TutorSource#1\endcsname\par}%
      \fi
    \fi
  \fi
}
\newcommand{\BankTask}[3]{%
  \refstepcounter{banktask}\par\medskip
  \sbox{\BankTaskBox}{\begin{minipage}{\linewidth}%
    \BankTaskContent{#1}{#2}{#3}\end{minipage}}%
  \ifdim\dimexpr\ht\BankTaskBox+\dp\BankTaskBox\relax>\dimexpr\textheight-20mm\relax
    \Needspace{35mm}\BankTaskContent{#1}{#2}{#3}%
  \else
    \Needspace{\dimexpr\ht\BankTaskBox+\dp\BankTaskBox\relax}%
    \noindent\usebox{\BankTaskBox}\par
  \fi
  \medskip
}
\begin{document}
\BookSubsection{1.2.1}{Закон равномерного прямолинейного движения}


% Source sdamgia-32279: https://phys-ege.sdamgia.ru/problem?id=32279
\BankTask{sdamgia-32279}{
Координата тела меняется с течением времени согласно закону $x = 4 - 2 t,$ где все величины выражены в СИ. Определите проекцию скорости $\upsilon _{x}$ этого тела. (В метрах в секунду.)
}{
\smallskip\noindent\textbf{Решение}\par
Зависимость $x(t)$ линейная, следовательно, тело движется равномерно и прямолинейно. Уравнение такого движения имеет вид $x ( t ) =x_0 + v_x t.$ Делаем вывод, что $v _x= - 2\,\text{м/с}.$

Ответ: $-2$.
}

% Source sdamgia-3544: https://phys-ege.sdamgia.ru/problem?id=3544
\BankTask{sdamgia-3544}{
Движение двух велосипедистов задано уравнениями $x_1 = 2t ( \text{м} )$ и $x_2 = 100 - 8t ( \text{м} ) .$ Найдите координату $x$ места встречи велосипедистов. Велосипедисты двигаются вдоль одной прямой. (В метрах.)
}{
\smallskip\noindent\textbf{Решение}\par
Встреча двух велосипедистов означает, что у них в некоторый момент времени совпадут координаты. Определим, когда именно произойдет встреча, для этого решим уравнение $x_1 ( t ) =x_2 ( t ) \Longleftrightarrow  2t=100 - 8t \Longleftrightarrow  t=10.$ Теперь не составляет труда определить координату места встречи: 

\[x_{\text{встреча}}=x_1 ( 10 ) =2 \cdot  10 ( \text{м} ) =20\,\text{м}=x_2 ( 10 ) =100 - 8 \cdot  10=20\,\text{м}.\]

Ответ: $20$.
}

% Source sdamgia-3354: https://phys-ege.sdamgia.ru/problem?id=3354
\BankTask{sdamgia-3354}{
Четыре тела двигались по оси Ох. В таблице представлена зависимость их координат от времени.

\begin{center}\small\renewcommand{\arraystretch}{1.25}\begin{tabular}{|c|c|c|c|c|c|c|}\hline
$t$, с & $0$ & $1$ & $2$ & $3$ & $4$ & $5$\\\hline
$x_{1}$, м. & $0$ & $2$ & $4$ & $6$ & $8$ & $10$\\\hline
$x_{2}$, м & $0$ & $0$ & $0$ & $0$ & $0$ & $0$\\\hline
$x_{3}$, м & $0$ & $1$ & $4$ & $9$ & $16$ & $25$\\\hline
$x_{4}$, м & $0$ & $2$ & $0$ & -$2$ & $0$ & $2$\\\hline
\end{tabular}\end{center}

У какого из тел скорость могла быть постоянна и отлична от нуля?
}{
\smallskip\noindent\textbf{Решение}\par
Таблица содержит информацию о положениях тел только в отдельные моменты времени. В промежутках между указанными временами тела могли двигаться абсолютно произвольно. Определим тело, скорость которого могла бы быть постоянной и отличной от нуля. При движении с постоянной скоростью координата тела за равные промежутки времени изменяется одинаково. Из таблицы видно, что этому свойству удовлетворяет только первое тело. Значит, скорость первого тела могла быть постоянной и не равной нулю.

 

Ответ: $1$.
}

% Source sdamgia-110: https://phys-ege.sdamgia.ru/problem?id=110
\BankTask{sdamgia-110}{
\begin{center}\begin{tikzpicture}[task graph,x=0.8cm,y=0.8cm]\draw[task grid,xstep=1,ystep=1](0,0)grid(6,4);\draw[task axis](0,0)--(6.45,0)node[right]{$t,\,\text{с}$};\draw[task axis](0,0)--(0,4.45)node[above]{$S,\,\text{м}$};\node[below left]at(0,0){$0$};\fill(1,0)circle(1pt);\node[below]at(1,0){$1$};\fill(2,0)circle(1pt);\node[below]at(2,0){$2$};\fill(3,0)circle(1pt);\node[below]at(3,0){$3$};\fill(0,1)circle(1pt);\node[left]at(0,1){$1$};\fill(0,2)circle(1pt);\node[left]at(0,2){$2$};\draw[task curve](0,0)--(1,1)--(3,1)--(5,2);\end{tikzpicture}\end{center}На рисунке представлен график зависимости пути от времени. Определите по графику скорость движения велосипедиста в интервале от момента времени $1\,\text{с}$ до момента времени $3\,\text{с}$ после начала движения. (В метрах в секунду.)
}{
\smallskip\noindent\textbf{Решение}\par
Из графика видно, что в интервале от момента времени $1\,\text{с}$ до момента времени $3\,\text{с}$ после начала движения путь велосипедиста не изменялся. Следовательно, на этом интервале времени велосипедист не двигался, его скорость была равна нулю.

Ответ: $0$.
}

% Source sdamgia-3867: https://phys-ege.sdamgia.ru/problem?id=3867
\BankTask{sdamgia-3867}{
Тела $1$ и $2$ двигаются вдоль оси $x$. На рисунке изображены графики зависимости координат движущихся тел $1$ и $2$ от времени $t$. 

\begin{center}\begin{tikzpicture}[task graph,x=0.45cm,y=0.025cm]\draw[task grid,xstep=1,ystep=20](0,0)grid(12,180);\draw[task axis](0,0)--(12.45,0)node[right]{$t,\,\text{с}$};\draw[task axis](0,0)--(0,189.0)node[above]{$x,\,\text{м}$};\node[below left]at(0,0){$0$};\fill(5,0)circle(1pt);\node[below]at(5,0){$5$};\fill(10,0)circle(1pt);\node[below]at(10,0){$10$};\fill(0,20)circle(1pt);\node[left]at(0,20){$20$};\fill(0,40)circle(1pt);\node[left]at(0,40){$40$};\fill(0,60)circle(1pt);\node[left]at(0,60){$60$};\fill(0,80)circle(1pt);\node[left]at(0,80){$80$};\fill(0,100)circle(1pt);\node[left]at(0,100){$100$};\fill(0,120)circle(1pt);\node[left]at(0,120){$120$};\fill(0,140)circle(1pt);\node[left]at(0,140){$140$};\draw[task curve](0,140)--(10,0);\draw[task curve](0,0)--(10,40);\node at(1.8,123){$1$};\node at(4,29){$2$};\end{tikzpicture}\end{center}

Чему равен модуль скорости $1$ относительно тела $2$? (В метрах в секунду.)
}{
\smallskip\noindent\textbf{Решение}\par
Используя график, определим проекции скоростей обоих тел. Для тела $1$ имеем: 

\[v_{x1}= \frac{\Delta x}{\Delta t} = \frac{x ( 10\,\text{с} ) - x ( 0\,\text{с} )}{10\,\text{с} - 0\,\text{с}} = \frac{0\,\text{м} - 140\,\text{м}}{10\,\text{с}} = - 14\,\text{м/с}.\]

Для тела $2$: 

\[v_{x2}= \frac{\Delta x}{\Delta t} = \frac{x ( 10\,\text{с} ) - x ( 0\,\text{с} )}{10\,\text{с} - 0\,\text{с}} = \frac{40\,\text{м} - 0\,\text{м}}{10\,\text{с}} =4\,\text{м/с}.\]

Таким образом, модуль скорости одного тела относительно другого равен:

\[v _{\text{отн}}=|\vec v _2 - \vec v _1|=| v_{x2} - v_{x1}|=|4\,\text{м/с} - ( - 14\,\text{м/с} ) |=18\,\text{м/с}.\]

Ответ: $18$.
}

% Source sdamgia-7995: https://phys-ege.sdamgia.ru/problem?id=7995
\BankTask{sdamgia-7995}{
\begin{center}\begin{tikzpicture}[task graph,x=0.065cm,y=0.012cm]\draw[task grid,xstep=10,ystep=50](0,0)grid(90,350);\draw[task axis](0,0)--(94.5,0)node[right]{$t,\,\text{с}$};\draw[task axis](0,0)--(0,372.5)node[above]{$S,\,\text{м}$};\node[below left]at(0,0){$0$};\fill(10,0)circle(1pt);\node[below]at(10,0){$10$};\fill(20,0)circle(1pt);\node[below]at(20,0){$20$};\fill(30,0)circle(1pt);\node[below]at(30,0){$30$};\fill(40,0)circle(1pt);\node[below]at(40,0){$40$};\fill(50,0)circle(1pt);\node[below]at(50,0){$50$};\fill(60,0)circle(1pt);\node[below]at(60,0){$60$};\fill(70,0)circle(1pt);\node[below]at(70,0){$70$};\fill(0,50)circle(1pt);\node[left]at(0,50){$50$};\fill(0,100)circle(1pt);\node[left]at(0,100){$100$};\fill(0,150)circle(1pt);\node[left]at(0,150){$150$};\fill(0,200)circle(1pt);\node[left]at(0,200){$200$};\fill(0,250)circle(1pt);\node[left]at(0,250){$250$};\draw[task curve](0,0)--(10,50)--(30,50)--(50,100)--(70,250);\end{tikzpicture}\end{center}На рисунке представлен график зависимости пути $S$ велосипедиста от времени $t$. Найдите скорость велосипедиста в интервале времени от $50$ до $70\,\text{с}$.
}{
\smallskip\noindent\textbf{Решение}\par
За время от $50$ до $70\,\text{с}$ велосипедист проехал $250 - 100 = 150\,\text{м}$, значит, его скорость равна $150\,\text{м}:20\,\text{с}=7{,}5\,\text{м/с}$.

Ответ: $7{,}5$.
}

% Source sdamgia-3783: https://phys-ege.sdamgia.ru/problem?id=3783
\BankTask{sdamgia-3783}{
\begin{center}\begin{tikzpicture}[task graph,x=0.5cm,y=0.02cm]\draw[task grid,xstep=2,ystep=40](0,0)grid(12,240);\draw[task axis](0,0)--(12.9,0)node[right]{$t,\,\text{с}$};\draw[task axis](0,0)--(0,258.0)node[above]{$S,\,\text{м}$};\node[below left]at(0,0){$0$};\fill(2,0)circle(1pt);\node[below]at(2,0){$2$};\fill(4,0)circle(1pt);\node[below]at(4,0){$4$};\fill(6,0)circle(1pt);\node[below]at(6,0){$6$};\fill(8,0)circle(1pt);\node[below]at(8,0){$8$};\fill(0,80)circle(1pt);\node[left]at(0,80){$80$};\fill(0,160)circle(1pt);\node[left]at(0,160){$160$};\draw[task curve](0,0)--(8,160);\draw[task curve](0,0)--(6.4,192);\node at(7,130){$1$};\node at(5.2,181){$2$};\end{tikzpicture}\end{center}На рисунке представлены графики зависимости пройденного пути от времени для двух тел. На какую величину $\Delta $$\upsilon $ скорость второго тела $\upsilon _{2}$ больше скорости первого тела $\upsilon _{1}$? (В метрах в секунду.)
}{
\smallskip\noindent\textbf{Решение}\par
Из графика видно, что для обоих тел пройденный путь линейно зависит от времени, а значит, оба тела двигались с постоянными по величине скоростями. Модуль скорости первого тела равен $v _1= \frac{160\,\text{м}}{8\,\text{с}} =20\,\text{м/с}.$ Скорость же второго тела: $v _2= \frac{120\,\text{м}}{4\,\text{с}} =30\,\text{м/с}.$ Следовательно, скорость второго тела больше скорости первого тела на величину $\Delta v = v _2 - v _1=30\,\text{м/с} - 20\,\text{м/с} =10\,\text{м/с}.$ 

Ответ: $10$.
}

% Source sdamgia-11807: https://phys-ege.sdamgia.ru/problem?id=11807
\BankTask{sdamgia-11807}{
Расстояние между пунктами А и В равно $30\,\text{км}$. Из пункта А в направлении пункта В выезжает мотоциклист со скоростью $50\,\text{км/ч}$. Одновременно из пункта В в том же направлении, что и мотоциклист, выезжает трактор со скоростью $20\,\text{км/ч}$. На каком расстоянии от пункта А мотоциклист догонит трактор? (В километрах.)
}{
\smallskip\noindent\textbf{Решение}\par
Для того чтобы догнать трактор, мотоциклисту нужно сократить расстояние в $30\,\text{км}$. Относительная скорость мотоциклиста относительно трактора равна $50\,\text{км/ч} - 20\,\text{км/ч} = 30\,\text{км/ч}$. Следовательно, мотоциклист догонит трактор за $\frac{30\,\text{км}}{30\,\text{км/ч}} = 1\,\text{ч}.$ За это время мотоциклист проедет $50\,\text{км}$.

Ответ: $50$.
}

% Source sdamgia-9752: https://phys-ege.sdamgia.ru/problem?id=9752
\BankTask{sdamgia-9752}{
Скорость течения широкой реки $3{,}6\,\text{км/ч}$. Под каким углом к направлению течения реки лодочник должен направлять лодку, скорость которой относительно воды равна $2\,\text{м/с}$, чтобы за $15$ минут ее снесло по направлению течения на $1{,}8\,\text{км}$?
}{
\smallskip\noindent\textbf{Решение}\par
При движении под углом $\alpha $ к течению реки, лодка будет двигаться по реке со скоростью 

\[V=V_{\text{теч}}+V_{\text{л}}\cos\alpha=3{,}6\cdot\frac{1000}{3600}+2\cos\alpha=(1+2\cos\alpha)\,\text{м/с}.\]

Тогда за время $t$ лодка переместится на расстояние

\[S=Vt= ( 1 + 2 \cdot  \cos  \alpha  ) t,\]

откуда

\[\cos  \alpha  = \frac{\frac{S}{t} - 1}{2} = \frac{1}{2} ( \frac{1800}{900} - 1 ) = \frac{1}{2} .\]

Лодочник должен направлять лодку под углом $60^\circ$ к направлению реки.

Ответ: $60$.
}

% Source sdamgia-23251: https://phys-ege.sdamgia.ru/problem?id=23251
\BankTask{sdamgia-23251}{
Школьник, возвращаясь домой с занятий по подготовке к ЕГЭ по физике, сел на вокзале в стоящую электричку. В это время пошел сильный снег, и вдоль поезда начал дуть ветер с постоянной скоростью. При этом поток падающих за окном снежинок выглядел так, как показано на рис. $1$. Потом поезд поехал, и после его разгона наблюдаемая через окно картина изменилась (см. рис. $2$). Углы наклона к горизонту прямолинейных траекторий снежинок в первом и во втором случаях были равны $\alpha  = 45 ^\circ$ и $\beta  = 10 ^\circ.$ С какой скоростью ехал поезд, если снежинки падают в неподвижном воздухе со скоростью $V = 1\,\text{м/с}$, электричка едет прямолинейно, а скорость и направление ветра не изменяются?
\begin{center}\begin{tikzpicture}[>=stealth,line width=.5pt]

\foreach\shift/\angle/\label in{0/45/1,5/10/2}{\begin{scope}[xshift=\shift cm]\draw(0,0)rectangle(4,3);\foreach\x in{.5,1.2,1.9,2.6,3.3}\foreach\y in{.4,1,1.6,2.2,2.8}\draw(\x,\y)--++(\angle:.15);\draw[thin](1.1,1.5)--(3,1.5);\draw[->,thick](2.6,1.5)--++(180+\angle:1.2);\draw(2.15,1.5)arc(180:180+\angle:.45);\node[below]at(2,-.15){Рис.\ \label};\end{scope}}
\node at(1.9,1.25){$\alpha$};\node at(6.9,1.23){$\beta$};
\end{tikzpicture}\end{center}

}{
\smallskip\noindent\textbf{Решение}\par
$1$. В неподвижной системе отсчета, связанной со стоящим поездом (рис. $1$), скорость снежинок $\overrightarrow V$ по закону сложения скоростей в кинематике складывается из двух составляющих — скорости $\overrightarrow V$ падения снежинки по вертикали и горизонтальной скорости ветра $\vec V_{\text{в}}:\quad\vec V_1=\vec V+\vec V_{\text{в}}.$ Поскольку векторы $\overrightarrow V$ и $\overrightarrow V_{\text{в}}$ взаимно перпендикулярны, а суммарный вектор $\overrightarrow V_1$ наклонен под углом $\alpha  = 45 ^\circ$ к горизонту, то модули $\overrightarrow V$ и $\overrightarrow V_{\text{в}}$ одинаковы и равны $V_{\text{в}}$ = $V = 1\,\text{м/с}$ — ветер дует влево с горизонтальной скоростью $1\,\text{м/с}$.

$2$. В системе отсчета, движущейся вместе с поездом (рис. $2$), снежинки приобретают относительно него дополнительную горизонтальную скорость $V_{\text{п}}$, направленную также влево (против вектора скорости поезда $\overrightarrow V_{\text{п}}$) и складывающуюся со скоростью ветра.

$3$. При этом вертикальная составляющая скорости снежинок остается равной $V = 1\,\text{м/с}$. Таким образом, как следует из прямоугольного треугольника для сложения горизонтальной и вертикальной составляющих скорости снежинки относительно поезда, $V = ( V_{\text{п}} + V_{\text{в}} ) \cdot  \tan  \beta  ,$ откуда $V_{\text{п}}=V/\tan\beta-V_{\text{в}}\approx1/0{,}176-1\approx4{,}67\,\text{м/с}\approx17\,\text{км/ч}.$

Ответ: $V_{\text{п}}\approx17\,\text{км/ч}.$
}


% Source B51AF2: https://ege.fipi.ru/bank/questions.php?search=1&proj=BA1F39653304A5B041B656915DC36B38&qid=B51AF2
\BankTask{B51AF2}{
На рисунке приведён график зависимости координаты тела $x$ от времени $t$ при его прямолинейном движении вдоль оси $Ox$. Определите проекцию скорости этого тела на ось $Ox$ в промежутке времени от $4$ до $6\,\text{с}$, в $\text{м}/\text{с}$.
\begin{center}
\begin{tikzpicture}[x=0.35cm,y=0.14cm,>=stealth]
\draw[step=1,black!85,thin,xstep=2,ystep=5] (0,-15) grid (14,15);
\draw[->] (0,0)--(16,0) node[right] {$t,\,\text{с}$};
\draw[->] (0,-16)--(0,20) node[above] {$x,\,\text{м}$};
\foreach \t in {2,4,6,8,10,12,14} \node[below] at (\t,0) {$\t$};
\foreach \h in {-15,-10,-5,0,5,10,15} \node[left] at (0,\h) {$\h$};
\draw[very thick] (0,-10)--(4,-10)--(8,10)--(14,-5);
\end{tikzpicture}
\end{center}
}{
\smallskip\noindent\textbf{Решение от GPT}\par
На участке от $4$ до $8\,\text{с}$ скорость постоянна. По графику $x(4)=-10\,\text{м}$, $x(6)=0$. Поэтому $v_x=\dfrac{\Delta x}{\Delta t}=\dfrac{0-(-10)}{6-4}=5\,\text{м}/\text{с}$.
}

\end{document}
